% ============================================================
% 1.2 Historical Background
% ============================================================
\subsection{Historical Background}\label{sub:HistoricalEstimates}

We already know how to solve the equations smoothly on a local time interval.
The open problem begins when we try to keep restarting that solution forever.
To restart it, we need some quantity the equations can stop from escaping.
The first quantity they control for every finite time is kinetic energy.

\divider{Leray's Energy Estimate}

Take the inner product of the momentum equation with \(u\) and integrate over space.
Nonlinear transport can move kinetic energy from one part of the flow to another, but the divergence-free condition makes its total contribution vanish.
Pressure can redirect the motion, but it cannot change the total energy either.
Viscosity is the term left over, and it removes energy through spatial variation.
For a smooth solution this calculation is an equality, so in particular
\begin{equation}\label{eq:LerayEnergyInequality}
  \lVert u(t)\rVert_{L^2}^{2}
  +2\nu\int_0^t\lVert\nabla u(s)\rVert_{L^2}^{2}\,ds
  \leq
  \lVert u_0\rVert_{L^2}^{2}.
\end{equation}
The first term keeps the kinetic energy bounded at each time.
The second pays for the total spatial variation dissipated by viscosity over the whole interval.

Leray then carried this inequality through a sequence of smooth approximations that only converge weakly.
Weak lower semicontinuity keeps the left side from being lost in the limit, which is enough to construct a finite-energy weak solution on \(\mathbb R^3\) for all time \parencite{Leray1934}.
On every finite interval this gives
\[
  u\in L^\infty(0,T;L^2),
  \qquad
  \nabla u\in L^2(0,T;L^2).
\]
Leray reached every finite time, but only in the energy class.
The estimate does not keep the solution classically differentiable, and it does not give uniqueness in three dimensions.
Its dissipation already points at the first derivative that can still escape.

\divider{The Enstrophy Problem}

For a divergence-free decaying field,
\[
  \lVert\nabla u\rVert_{L^2}^{2}
  =
  \lVert\omega\rVert_{L^2}^{2},
  \qquad
  \omega:=\nabla\times u.
\]
The squared \(L^2\) size of the vorticity is the enstrophy.
Leray's estimate controls the area under its time graph.
It does not control the height of that graph at each time, so very large enstrophy on very short intervals is still compatible with the energy inequality.

Taking the curl of the momentum equation shows where that growth can come from:
\begin{equation}\label{eq:HistoricalVorticityEquation}
  \partial_t\omega+(u\cdot\nabla)\omega
  =
  (\omega\cdot\nabla)u+\nu\Delta\omega.
\end{equation}
Transport carries the vorticity along the flow, while viscosity diffuses it.
The extra three-dimensional term stretches it.
In two dimensions this stretching term is absent, and the enstrophy estimate closes.
In three dimensions a surrounding gradient of \(u\) can elongate a vortex line and increase its intensity while viscosity is trying to smooth it.
At the level of total enstrophy,
\[
  \frac12\frac{d}{dt}\lVert\omega\rVert_{L^2}^{2}
  +\nu\lVert\nabla\omega\rVert_{L^2}^{2}
  =
  \int_{\mathbb R^3}(\omega\cdot\nabla)u\cdot\omega\,dx.
\]
Viscosity contributes a definite loss, while vortex stretching can either raise or lower the enstrophy.
The global energy bound does not compare the two strongly enough to keep the enstrophy bounded at every time.

\divider{Higher Sobolev Continuation}

Local theory tells us what would rule out this escape.
Suppose a smooth solution exists on \([0,T)\) and, for some \(s>5/2\),
\begin{equation}\label{eq:HistoricalSobolevContinuation}
  \sup_{0\leq t<T}\lVert u(t)\rVert_{H^s}<\infty.
\end{equation}
Sobolev embedding then controls \(\nabla u\), the differentiated energy estimates stay finite, and the local existence theorem can be run again at time \(T\).
Classical continuation theory gives the same restart under the accumulated condition \parencite{FujitaKato1964,Kato1984}
\[
  \int_0^T\lVert\nabla u(t)\rVert_{L^\infty}\,dt<\infty.
\]

If either bound stays finite, then \(T\) is not the end of the smooth solution.
So if \(T_*\) is finite and maximal, neither restart quantity can remain finite as \(t\uparrow T_*\).
What is still missing is a way to derive either bound from arbitrary smooth initial data for all time.

\divider{Caffarelli\textendash Kohn\textendash Nirenberg Partial Regularity}

Without a global restart bound, Caffarelli, Kohn, and Nirenberg moved the energy argument into a shrinking neighbourhood of a spacetime point.
For \(z_0=(x_0,t_0)\), write
\[
  Q_r(z_0)
  :=
  B_r(x_0)\times(t_0-r^2,t_0).
\]
Suitable weak solutions satisfy a local energy inequality on these parabolic cylinders.
Their epsilon-regularity argument says that if the scale-invariant quantities built from \(u\) and \(p\) are small enough on \(Q_r(z_0)\), then the solution is regular on a smaller cylinder.

Now read the same statement from a singular point.
Every sufficiently small cylinder around that point must retain a definite amount of scale-invariant concentration; if even one cylinder fell below the threshold, the point would be regular.
A covering argument then confines the possible singular set to a set of zero one-dimensional parabolic Hausdorff measure \parencite{CaffarelliKohnNirenberg1982}.
The singular set can be this sparse and still contain one point.
That one point is enough to break global smoothness, and it would have to sustain concentration through every smaller scale.
Ruling out that last point requires a bound that controls its spatial intensity and temporal persistence together.

\divider{Ladyzhenskaya\textendash Prodi\textendash Serrin Criteria}

The Ladyzhenskaya\textendash Prodi\textendash Serrin estimates measure those two demands together.
For \(q>3\), H\"older's inequality and interpolation bound the nonlinear term by
\[
  C\lVert u\rVert_{L^q}
  \lVert\nabla u\rVert_{L^2}^{1-3/q}
  \lVert\Delta u\rVert_{L^2}^{1+3/q}.
\]
Because \(1+3/q<2\), Young's inequality absorbs the highest derivative into the viscous term and leaves
\[
  \frac{d}{dt}\lVert\nabla u\rVert_{L^2}^{2}
  +\nu\lVert\Delta u\rVert_{L^2}^{2}
  \leq
  C_{\nu,q}\lVert u\rVert_{L^q}^{\frac{2q}{q-3}}
  \lVert\nabla u\rVert_{L^2}^{2}.
\]
The exponent \(2q/(q-3)\) tells us how much time-integrability of \(\lVert u\rVert_{L^q}\) is needed: the coefficient is integrable when
\begin{equation}\label{eq:HistoricalLPSCriterion}
  u\in L^p(0,T;L^q(\mathbb R^3)),
  \qquad
  \frac{2}{p}+\frac{3}{q}\leq1,
  \qquad
  q>3.
\end{equation}
Gr\"onwall's inequality then keeps the enstrophy finite, and the continuation theory carries the solution through \(T\).
This gives the Ladyzhenskaya\textendash Prodi\textendash Serrin regularity criterion for Leray\textendash Hopf weak solutions \parencite{Prodi1959,Serrin1963,Ladyzhenskaya1967}.

The exponent \(q\) measures how sharply \(u\) can concentrate in space, while \(p\) measures how long that concentration can persist.
At equality, their mixed norm is unchanged by the Navier\textendash Stokes scaling.
Any finite singular endpoint must make every admissible Serrin mixed norm infinite.
But \(2q/(q-3)\) blows up at \(q=3\), so the estimate fails exactly where the critical spatial norm begins.

\divider{The Escauriaza\textendash Seregin\textendash \v{S}ver{\'a}k Endpoint}

Under the Navier\textendash Stokes rescaling
\begin{equation}\label{eq:HistoricalNavierStokesScaling}
  u_\lambda(x,t)
  :=
  \lambda u(\lambda x,\lambda^2t),
\end{equation}
the spatial \(L^3\) norm does not change:
\[
  \lVert u_\lambda(t)\rVert_{L^3}
  =
  \lVert u(\lambda^2t)\rVert_{L^3}.
\]
A concentration can be magnified back to unit scale without making this norm any smaller.

Now assume that \(T\) is a first singular time while the \(L^3\) norm stays bounded up to \(T\).
Magnify the solution around the suspected concentration.
The critical bound survives every magnification, so compactness produces an ancient limiting solution with the same bound.
The concentration used to choose the scales makes that limit nonzero.
The limiting field vanishes at its final time, so its vorticity does too.
Backward uniqueness carries that vanishing backward outside a large ball, and unique continuation forces the vorticity to vanish everywhere.
The limiting field is then zero.
The limit cannot be both nonzero and zero.
Escauriaza, Seregin, and \v{S}ver{\'a}k used this contradiction to prove the critical endpoint criterion \parencite{EscauriazaSereginSverak2003}: a Leray\textendash Hopf weak solution on \(\mathbb R^3\) remains regular through \(T\) whenever
\[
  u\in L^\infty(0,T;L^3(\mathbb R^3)).
\]

The endpoint result says the critical norm must become unbounded near any finite maximal time.
That still leaves open the possibility that it rises along isolated times and repeatedly falls back below some fixed level.
Seregin ruled out that behaviour for smooth finite-energy solutions on \(\mathbb R^3\) \parencite{Seregin2012}: if \(T_*\) is finite and maximal, then
\begin{equation}\label{eq:SereginLThreeDivergence}
  \lim_{t\uparrow T_*}\lVert u(t)\rVert_{L^3(\mathbb R^3)}
  =\infty.
\end{equation}
So the critical concentration has to leave every finite level and stay above it as the endpoint approaches.

Seregin's result gives no rate for this divergence.
Tao proved that for infinitely many times approaching any finite maximal time, the \(L^3\) norm has to exceed a positive power of a triple logarithm of the inverse time remaining \parencite{Tao2021Quantitative}.

\divider{Where the Problem Sits Today}

Now compare the critical norm with the global quantity Leray controls.
At corresponding times,
\[
  \lVert u_\lambda(t)\rVert_{L^3}
  =
  \lVert u(\lambda^2t)\rVert_{L^3},
  \qquad
  \lVert u_\lambda(t)\rVert_{L^2}^{2}
  =
  \lambda^{-1}\lVert u(\lambda^2t)\rVert_{L^2}^{2}.
\]
Magnifying a smaller scale preserves the critical concentration while the rescaled kinetic energy becomes smaller.
Thus the global \(L^2\) bound alone cannot put a scale-independent ceiling on the critical \(L^3\) concentration.

Smaller structures have larger gradients and stronger viscous dissipation, but nonlinear transport can keep transferring activity into those scales.
The unresolved question is whether viscous dissipation controls that transfer before a restart quantity diverges.

A frequency-doubling picture makes the timing visible, though it is only a heuristic.
If the active frequency doubles at each stage, its spatial scale halves and its natural parabolic time scale quarters.
The model times then form a convergent geometric series, so time alone does not rule out an unlimited sequence of scale changes.
That finite accumulation is the Zeno cascade picture.

The open problem is to use the original datum, viscosity, and equations to obtain a restart bound that stays finite through this possible scale accumulation.
